% 本文件由 LaTeX 排版 Agent 根据有效 translation.json 自动生成。
% author=Justin Martyr; work_id=0132; title=圣犹斯定殉道者残篇

\chapter{圣犹斯定殉道者残篇}

% structure: p0001 source=h1 target=omitted reason=page-title-already-rendered-by-parent

% structure: p0002 source=h2 target=section reason=relative-heading-rank; base=h2

\section{一}

最令人钦佩的\Person{犹斯定}正确地说，上述的魔鬼如同强盗。——\Person{塔齐安}的\WorkTitle{致希腊人书}第十八章。\par

% structure: p0004 source=h2 target=section reason=relative-heading-rank; base=h2

\section{二}

犹斯定在他的\WorkTitle{驳斥马尔基昂}一书中说得对：如果主自己宣布了另一位神，而不是那位塑造者和造物主（世界的），以及我们的养育者，他就不会相信主。但是，既然从这独一的天主——他创造了这个世界并形成了我们，涵容并管理万有——那里，独生子来到我们这里，在他自己内总括了他的工程，那么我对他的信德是坚定的，我对圣父的爱德是不动摇的，天主将这两者都赐予了我们。——\Person{伊里奈乌斯}\WorkTitle{异端}卷四第六章。\par

% structure: p0006 source=h2 target=section reason=relative-heading-rank; base=h2

\section{三}

犹斯定说得对：在主来临之前，撒殚从未敢亵渎天主，因为他尚不确定自己的定罪，因为关于他的这事只是由先知们在比喻和寓言中宣布的。但在主来临之后，撒殚从基督及其宗徒的讲话中清楚地知道，永远的火是为那自愿离开天主的人以及所有不悔改而坚持背道的人预备的，于是，他借着这样一个人，如同已经定罪一样，亵渎那施加审判于他的天主，并将他背道的罪归咎于他的造物主，而不是归咎于他自己的意志和偏好。——\Person{伊里奈乌斯}\WorkTitle{异端}卷五第二十六章。\par

% structure: p0008 source=h2 target=section reason=relative-heading-rank; base=h2

\section{四}

他在阐述魔鬼不休地谋害我们的原因时宣称：在主来临之前，魔鬼并不十分清楚他自己惩罚的尺度，因为先知们只是隐晦地宣布了它；例如依撒意亚，在以亚述人的人物中，悲惨地揭示了要对魔鬼施行的程序。但当主显现，魔鬼清楚地知道永远的火是为他和他的天使所存留和预备的，他便开始不停地谋划反对信友们，渴望在他的背道中有许多同伴，以免他独自承受定罪的耻辱，用这种冷酷而恶毒的安慰来安慰自己。——出自\Person{安提约基雅的若望}的著作。\par

% structure: p0010 source=h2 target=section reason=relative-heading-rank; base=h2

\section{五}

纳波利的犹斯定，一位在年龄和德行上都与宗徒们相去不远的人，说：可朽的继承，而不朽的则承受；肉身确实死去，但天国活着。——出自\Person{美多德}的\WorkTitle{论复活}，见于\Person{佛提乌}。\par

% structure: p0012 source=h2 target=section reason=relative-heading-rank; base=h2

\section{六}

在天主那里没有狭窄，也没有任何不是绝对完美的东西。——出自\Person{犹斯定}著作手稿。\par

% structure: p0014 source=h2 target=section reason=relative-heading-rank; base=h2

\section{七}

我们不会因不认识天主而伤害天主，但会剥夺自己对他的友谊。\par

% structure: p0016 source=h2 target=section reason=relative-heading-rank; base=h2

\section{八}

教师的不熟练对他的门徒是毁灭性的，门徒的粗心给教师带来危险，尤其是如果他们的疏忽归因于他的知识不足。\par

% structure: p0018 source=h2 target=section reason=relative-heading-rank; base=h2

\section{九}

灵魂很难被召回它已堕落的美好事物，也很难被拖离它已习惯的邪恶。如果你任何时候表现出自责的倾向，那么也许通过忏悔的良药，我会对你抱有良好的希望。但当你完全蔑视恐惧，并轻蔑地拒绝基督本身的信仰时，你倒不如从未从母腹中出生。——出自\Person{达玛森的若望}的著作。\par

% structure: p0020 source=h2 target=section reason=relative-heading-rank; base=h2

\section{十}

那两只鸟表徵基督，既是死的人，又是活的天主。他被比作鸟，因为人被理解为并宣告为来自高处和天上。那活的鸟，被蘸在死鸟的血中，之后被释放了。因为那活的、神圣的圣言是在被钉死且死亡的圣殿（肉身）中，作为苦难的分享者，而就天主而言，却是不受苦难的。\par

在流水中发生的事——其中木头、牛膝草和朱红色线被蘸——预表了基督在十字架上为那些被圣神、水和血所洒的人的救恩而流的血之苦难。因此，为洁净提供的材料主要不是针对癞病，而是针对罪的赦免，以便癞病被理解为罪的象征，而被献祭的东西则是那为罪而被献祭者的象征。\par

因此，他命令同时将朱红色线蘸在水中，从而预言肉身将不再拥有其本性的（邪恶）属性。因此，也有那两只鸟，一只是在水中被献祭的，另一只是既蘸在血中又蘸在水中然后被放走的，正如关于公山羊的叙述也是如此。\par

那只被放走的公山羊预表那除去人类罪恶的一位。但这两只公山羊包含了一个表徵：天主降生成人的同一救世计划。因为他为我们的悖逆而被刺伤，他担当了多人的罪，他为我们的过犯而被交付。——出自\Person{犹斯定}著作手稿。\par

% structure: p0025 source=h2 target=section reason=relative-heading-rank; base=h2

\section{十一}

天主起初造人时，将本性之事悬于其意志之上，并藉一条诫命进行试验。因为祂规定：人若遵守此诫，便可分享不死之存在；但若违犯，则承受相反的命运。人既如此被造，立刻趋向违命，自然便成了朽坏的奴隶。朽坏既然成为本性中的固有之物，那愿意拯救人的一位，就必须是摧毁朽坏之因由者。这事除非藉着按本性而有的生命与那已承受朽坏者结合，从而摧毁朽坏，同时使那承受朽坏者将来永保不朽，否则别无他法。因此，圣言必须取得身体，为将我们从本性的朽坏之死中拯救出来。因为，如你所说，若祂仅凭点头就为我们挡开死亡，那么死亡固然因祂意志的表达而不会接近我们，但我们仍将再次成为可朽坏的，因为我们自身带有那本性的朽坏。——\Person{良提乌} \WorkTitle{驳尤提基派}，卷二。\par

% structure: p0027 source=h2 target=section reason=relative-heading-rank; base=h2

\section{12}

如同天主所造的一切身体固有影子一样，公义的天主也理应按照各人的功过，回报那些选择善事的人和那些偏爱恶事的人。——摘自\Person{大马士革的若望}著作。\par

% structure: p0029 source=h2 target=section reason=relative-heading-rank; base=h2

\section{13}

祂所说的并非在外邦土地上的外邦人，而是指那些附和外邦人的[百姓]，正如\BibleBook{耶肋米亚}{先知书}所说：\Quote{你离弃我，是件痛苦的事，主你的天主说，因为你从古时已折断你的轭，撕裂你的绳索，说：我不愿事奉你，却要登上每一座高丘，在每一棵树下行淫。}——摘自\Person{儒斯定}著作手稿。\par

% structure: p0031 source=h2 target=section reason=relative-heading-rank; base=h2

\section{14}

光明永不会成为黑暗，只要光明存在；关乎我们之事的真理也不会被质疑。因为真理是比一切都更强大的。凡能说真理却不说的人，必受天主的审判。——摘自\Person{大马士革的若望}手稿与著作。\par

% structure: p0033 source=h2 target=section reason=relative-heading-rank; base=h2

\section{15}

关于第七日，并未像其他日子一样说\Quote{有晚上，有早晨}，这明显是预指在它完成之前将要发生的圆满，正如教父们所宣告的，尤其是圣\Person{克莱孟}、\Person{依勒内}以及殉道者哲学家\Person{儒斯定}所论。他在极其智慧地诠释第六日的数目六时，断言人的理智灵魂及其五种感觉官能就是第六日的六项工程\OriginalTerm{六项工程}{six works}。由此，在详尽论述了数字六之后，他宣称天主所造的一切事物分为六类：即，理智且不朽之物，如天使；理性且有死之物，如人类；有感觉且无理性之物，如牲畜、鸟类和鱼类；能前进、运动且无感觉之物，如风、云、水、星辰；能生长且不动之物，如树木；以及无感觉且不动之物，如山岳、大地等。因为天主在天上和地上的一切受造物，都归属于这六类之一，并被它们所界定。——摘自\Person{阿纳斯塔西}著作。\par

% structure: p0035 source=h2 target=section reason=relative-heading-rank; base=h2

\section{16}

纯正的道理不能进入顽固悖逆的心；反而如同被击退，重新进入自身。\par

% structure: p0037 source=h2 target=section reason=relative-heading-rank; base=h2

\section{17}

正如身体的善是健康，灵魂的善就是知识，这知识确实是某种灵魂的健康，藉此达到与天主的相似。——摘自\Person{大马士革的若望}著作。\par

% structure: p0039 source=h2 target=section reason=relative-heading-rank; base=h2

\section{18}

屈服并顺从私欲是最卑贱的奴役，正如克制私欲是惟一的自由。\par

最大的善是无罪，其次是成义；但那些生活在不义中却长久不受惩罚的人，必被视为最不幸的。\par

套上轭的牲畜，若驭手无经验且恶劣，便不免被拖下悬崖；若他娴熟且卓越，它们便能得救。\par

哲学家的目标是尽可能相似于天主。——摘自\Person{安多尼·梅利萨}著作。\par

% structure: p0044 source=h2 target=section reason=relative-heading-rank; base=h2

\section{19}

[圣\Person{儒斯定}、哲学家及殉道者]之言，出自其\WorkTitle{护教书}第五部分：——诸君，我认为福乐不是别的，正是按真理生活。但除非我们认识事物的本性，我们就不能适当或按真理生活。\par

他们显然不明白：那因真信德从谬误进入真理的人，已真正认识了自己，并非如他们所说处于狂乱状态，而是摆脱了变幻无常的（并针对各种谬误）可变的朽坏，藉着单纯且永远同一的真理。——摘自\Person{大马士革的若望}著作。\par

\SourceRecord{Translated by Alexander Roberts.；Ante-Nicene Fathers；Vol. 1.；Edited by Alexander Roberts, James Donaldson, and A. Cleveland Coxe.；Buffalo, NY: Christian Literature Publishing Co.,；1885.；<http://www.newadvent.org/fathers/0132.htm>.}
